\documentclass[12pt,a4paper]{article}

\usepackage[utf8]{inputenc}
\usepackage[spanish]{babel}
\usepackage[T1]{fontenc}
\usepackage{geometry}
\usepackage{titlesec}
\usepackage{enumitem}
\usepackage{booktabs}
\usepackage{array}
\usepackage{xcolor}
\usepackage{hyperref}
\usepackage{parskip}
\usepackage{fancyhdr}
\usepackage{graphicx}

\geometry{
  top=3cm,
  bottom=2.5cm,
  left=3cm,
  right=2.5cm
}

\hypersetup{
  colorlinks=false
}

\titleformat{\section}
  {\large\bfseries}
  {\thesection.}{0.5em}{}[\titlerule]

\titleformat{\subsection}
  {\normalsize\bfseries}
  {\thesubsection.}{0.5em}{}

\pagestyle{fancy}
\fancyhf{}
\rhead{\small PropertyReservation}
\lhead{\small Seminario I}
\cfoot{\thepage}
\renewcommand{\headrulewidth}{0.4pt}

% --- DOCUMENTO ---

\begin{document}

% ---------------------------------------------------------
% CARÁTULA
% ---------------------------------------------------------
\begin{titlepage}
  \centering
  \vspace*{1cm}

  \includegraphics[width=2.5cm]{logo_ips.png}

  \vspace{0.5cm}

  {\large \textbf{Instituto Politécnico Superior ``Gral. San Martín''} \par}
  \vspace{0.2cm}
  {\normalsize Analista Universitario en Sistemas \par}
  \vspace{0.1cm}
  {\normalsize Seminario I \par}

  \vspace{1.5cm}

  {\LARGE \textbf{PropertyReservation} \par}
  \vspace{0.2cm}
  {\normalsize \textit{Descripción funcional del sistema} \par}

  \vspace{1.5cm}

  {\normalsize \textbf{Alumnos:} Matías Haspert y Virginia Bonservizi \par}

  \vspace{1.5cm}

  \tableofcontents

  \vfill
\end{titlepage}

\newpage

% ---------------------------------------------------------
% 1. DESCRIPCIÓN GENERAL
% ---------------------------------------------------------
\section{Descripción general del sistema}

\textbf{PropertyReservation} es un sistema de publicación y reserva de propiedades de alojamiento
temporario. Conecta a propietarios que quieren ofrecer sus inmuebles con huéspedes que buscan
dónde hospedarse en fechas específicas.

El sistema acompaña el proceso completo: el propietario publica su propiedad, el huésped la
encuentra y hace una reserva, luego sube un comprobante de pago que el dueño revisa, y al
finalizar la estadía el huésped puede dejar una reseña.

El sistema tiene dos interfaces: una \textbf{aplicación web} pensada para el uso general de
huéspedes y propietarios, y una \textbf{aplicación de escritorio} orientada a la gestión
administrativa y moderación de reseñas.

% ---------------------------------------------------------
% 2. ALCANCE
% ---------------------------------------------------------
\section{Alcance del sistema}

El sistema cubre los siguientes módulos:

\begin{itemize}[leftmargin=1.5em]
  \item \textbf{Gestión de usuarios:} registro e inicio de sesión con roles diferenciados.
  \item \textbf{Gestión de propiedades:} publicación, edición, carga de fotos y asignación de amenities.
  \item \textbf{Disponibilidad:} los propietarios configuran los períodos en que su propiedad acepta reservas.
  \item \textbf{Búsqueda y reserva:} los huéspedes buscan propiedades por destino y fechas, y generan reservas. Cada reserva pasa por estados sucesivos: pendiente de pago, comprobante enviado, confirmada, completada o cancelada.
  \item \textbf{Comprobantes de pago:} el huésped sube una imagen de comprobante; el propietario la aprueba o rechaza. El sistema no procesa pagos reales ni se integra con pasarelas de pago.
  \item \textbf{Reseñas:} al completarse una estadía, el huésped puede calificar la propiedad.
  \item \textbf{Panel de propietario:} resumen de propiedades y reservas pendientes de atención.
  \item \textbf{Aplicación de escritorio:} interfaz para el administrador, desde donde gestiona amenities, consulta estadísticas generales y modera reseñas.
\end{itemize}

\subsection{Fuera del alcance}

\begin{itemize}[leftmargin=1.5em]
  \item Integración con pasarelas de pago (tarjeta, PayPal, transferencias bancarias en tiempo real).
  \item Mensajería entre huéspedes y propietarios.
\end{itemize}

% ---------------------------------------------------------
% 3. ROLES Y FUNCIONALIDADES
% ---------------------------------------------------------
\section{Roles y funcionalidades}

El sistema tiene tres tipos de usuario. Un mismo usuario no puede ser huésped y propietario
al mismo tiempo; el rol se define en el momento del registro.

\subsection{Huésped}

\begin{itemize}[leftmargin=1.5em]
  \item \textbf{Registro e inicio de sesión}: crea una cuenta y accede con sus credenciales.
  \item \textbf{Búsqueda}: explora propiedades publicadas, filtra por destino y fechas, y ve los detalles de cada una (fotos, amenities, precio por noche y reseñas).
  \item \textbf{Reserva}: elige una propiedad con fechas disponibles y genera la reserva.
  \item \textbf{Comprobante de pago}: sube una imagen del comprobante para que el propietario la revise. No se realiza ningún pago a través del sistema.
  \item \textbf{Mis reservas}: consulta el estado de sus reservas actuales e históricas.
  \item \textbf{Reseñas}: al finalizar la estadía puede calificar la propiedad con 1 a 5 estrellas y un comentario.
\end{itemize}

\subsection{Propietario}

\begin{itemize}[leftmargin=1.5em]
  \item \textbf{Publicación de propiedades}: crea y edita sus propiedades con descripción, tipo, dirección, precio por noche y fotos.
  \item \textbf{Amenities}: indica qué servicios ofrece (Wi-Fi, estacionamiento, pileta, etc.) desde un listado predefinido.
  \item \textbf{Disponibilidad}: define los períodos en que cada propiedad acepta reservas.
  \item \textbf{Reservas recibidas}: ve las reservas de sus propiedades con los datos del huésped, fechas y estado.
  \item \textbf{Revisión de comprobantes}: aprueba o rechaza el comprobante que sube el huésped. Si lo aprueba, la reserva queda confirmada.
\end{itemize}

\subsection{Administrador}

Accede al sistema únicamente a través de la aplicación de escritorio.

\begin{itemize}[leftmargin=1.5em]
  \item \textbf{Estadísticas}: panel con totales de reservas, propiedades y usuarios registrados.
  \item \textbf{Amenities}: crea, edita y elimina las amenities disponibles en el sistema. No puede eliminar una comodidad que esté asignada a alguna propiedad.
  \item \textbf{Moderación de reseñas}: revisa y gestiona las reseñas publicadas en el sistema.
\end{itemize}

% ---------------------------------------------------------
% 4. TECNOLOGÍAS UTILIZADAS
% ---------------------------------------------------------
\section{Tecnologías utilizadas}

\begin{itemize}[leftmargin=1.5em]
  \item \textbf{React} (JavaScript) — interfaz web.
  \item \textbf{C\# con .NET 8} — servidor backend con API REST.
  \item \textbf{SQL Server} — base de datos relacional.
  \item \textbf{JWT} — autenticación y control de acceso por rol.
  \item \textbf{Windows Forms (C\#)} — aplicación de escritorio para el administrador.
\end{itemize}


\end{document}

